\section{Tác động dự kiến}

\subsection{Lợi ích xã hội/kinh doanh}

\begingroup
\footnotesize
\setlength{\tabcolsep}{3pt}
\renewcommand{\arraystretch}{1.25}

\begin{longtable}{|>{\raggedright\arraybackslash}p{0.23\textwidth}|
                  >{\raggedright\arraybackslash}p{0.43\textwidth}|
                  >{\raggedright\arraybackslash}p{0.26\textwidth}|}
\hline
\multicolumn{1}{|>{\centering\arraybackslash}p{0.23\textwidth}|}{\textbf{Nhóm hưởng lợi}} &
\multicolumn{1}{>{\centering\arraybackslash}p{0.43\textwidth}|}{\textbf{Lợi ích cụ thể}} &
\multicolumn{1}{>{\centering\arraybackslash}p{0.26\textwidth}|}{\textbf{Cách đo lường}} \\ \hline
\endfirsthead

\hline
\multicolumn{1}{|>{\centering\arraybackslash}p{0.23\textwidth}|}{\textbf{Nhóm hưởng lợi}} &
\multicolumn{1}{>{\centering\arraybackslash}p{0.43\textwidth}|}{\textbf{Lợi ích cụ thể}} &
\multicolumn{1}{>{\centering\arraybackslash}p{0.26\textwidth}|}{\textbf{Cách đo lường}} \\ \hline
\endhead

Sinh viên &
Biết mình đủ điều kiện cấp nào, còn thiếu gì, deadline nào đang mở; lưu thành tích từ sớm; tái sử dụng hồ sơ cấp cũ; được gợi ý hoạt động phù hợp để bù tiêu chí còn thiếu; có CV/portfolio từ minh chứng và danh hiệu đã được xác nhận &
Tỷ lệ hoàn thiện hồ sơ đúng hạn, thời gian phát hiện thiếu minh chứng, tỷ lệ sinh viên tham gia hoạt động được gợi ý, số CV/portfolio được xuất, mức hài lòng \\ \hline

Cán bộ cấp khoa/trường &
Giảm thời gian đọc minh chứng, lọc hồ sơ, tổng hợp danh sách &
Thời gian xử lý trung bình/hồ sơ, số hồ sơ xử lý/ngày, tỷ lệ hồ sơ bị trả bổ sung \\ \hline

Cán bộ cấp thành phố/tỉnh &
Nhận hồ sơ từ nhiều trường có cấu trúc chuẩn, dễ tổng hợp danh sách chính thức &
Thời gian chốt danh sách, số lỗi khi gửi lên cấp cao hơn \\ \hline

Cán bộ Trung ương &
Xem được lịch sử xét duyệt liên cấp, hồ sơ số hóa và danh sách chính thức &
Thời gian kiểm tra hồ sơ cấp Trung ương, tỷ lệ hồ sơ có đủ lịch sử/minh chứng \\ \hline

Đơn vị triển khai &
Chuẩn hóa quy trình nhưng vẫn linh hoạt theo điều kiện riêng &
Số đơn vị tham gia, số bộ điều kiện được cấu hình, tỷ lệ quy trình được số hóa \\ \hline

Hội Sinh viên/Đoàn - Hội phối hợp triển khai &
Có công cụ số hóa phong trào nhưng vẫn giữ quyền quản lý nghiệp vụ, chủ động cấu hình kế hoạch xét, nguồn hoạt động, bộ tiêu chí và báo cáo theo từng mùa &
Số đơn vị Hội Sinh viên tham gia thí điểm, số cán bộ được tập huấn, mức độ hài lòng của cán bộ, số góp ý nghiệp vụ được đưa vào phiên bản tiếp theo \\ \hline

Hệ thống phong trào Sinh viên 5 tốt &
Tăng minh bạch, giảm sai sót, tạo dữ liệu dài hạn về hành trình phấn đấu của sinh viên &
Số danh hiệu được ghi nhận theo cấp, số hồ sơ liên cấp, chất lượng báo cáo \\ \hline

\end{longtable}
\endgroup

\subsection{TAM - SAM - SOM hoặc người dùng tiềm năng}

\begingroup
\footnotesize
\setlength{\tabcolsep}{3pt}
\renewcommand{\arraystretch}{1.25}

\begin{longtable}{|>{\raggedright\arraybackslash}p{0.12\textwidth}|
                  >{\raggedright\arraybackslash}p{0.36\textwidth}|
                  >{\raggedright\arraybackslash}p{0.22\textwidth}|
                  >{\raggedright\arraybackslash}p{0.22\textwidth}|}
\hline
\multicolumn{1}{|>{\centering\arraybackslash}p{0.12\textwidth}|}{\textbf{Chỉ số}} &
\multicolumn{1}{>{\centering\arraybackslash}p{0.36\textwidth}|}{\textbf{Định nghĩa trong bài toán}} &
\multicolumn{1}{>{\centering\arraybackslash}p{0.22\textwidth}|}{\textbf{Ước tính sơ bộ}} &
\multicolumn{1}{>{\centering\arraybackslash}p{0.22\textwidth}|}{\textbf{Cơ sở ước tính}} \\ \hline
\endfirsthead

\hline
\multicolumn{1}{|>{\centering\arraybackslash}p{0.12\textwidth}|}{\textbf{Chỉ số}} &
\multicolumn{1}{>{\centering\arraybackslash}p{0.36\textwidth}|}{\textbf{Định nghĩa trong bài toán}} &
\multicolumn{1}{>{\centering\arraybackslash}p{0.22\textwidth}|}{\textbf{Ước tính sơ bộ}} &
\multicolumn{1}{>{\centering\arraybackslash}p{0.22\textwidth}|}{\textbf{Cơ sở ước tính}} \\ \hline
\endhead

TAM &
Toàn bộ sinh viên đại học/cao đẳng, cán bộ Đoàn - Hội và đơn vị triển khai phong trào Sinh viên 5 tốt trên toàn quốc &
2.000.000+ sinh viên; hàng trăm trường/đơn vị &
Phong trào có thể triển khai ở nhiều cấp: khoa, trường, thành phố/tỉnh, Trung ương \\ \hline

SAM &
Nhóm có thể tiếp cận trong 1-2 năm đầu: trường/khoa/tỉnh thành có nhu cầu số hóa xét duyệt &
100.000-300.000 sinh viên; 20-50 đơn vị &
Ưu tiên nơi đang xử lý bằng Google Form/Excel và có lượng hồ sơ lớn \\ \hline

SOM &
Phần có thể đạt sau cuộc thi: thí điểm tại 1-3 khoa/trường hoặc một cụm đơn vị nhỏ &
500-3.000 sinh viên; 300-1.000 hồ sơ/mùa xét &
Phù hợp kiểm chứng MVP 2 cấp khoa -{}> trường, OCR, dashboard và luồng liên cấp \\ \hline

\end{longtable}
\endgroup

\subsection{Ưu thế cạnh tranh}

\begin{itemize}
    \item \textbf{Bám sát quy trình thực tế:} 5-Star Eco xử lý đúng luồng khoa -{}> trường -{}> thành phố/tỉnh -{}> Trung ương, thay vì chỉ thu hồ sơ cuối kỳ.
    \item \textbf{Phân quyền theo cây tổ chức:} giải quyết bài toán mỗi cán bộ chỉ được xem hồ sơ thuộc đơn vị mình.
    \item \textbf{Tiêu chí linh hoạt nhưng có khung chung:} 5 nhóm tốt cố định, bộ điều kiện đạt cấu hình theo đơn vị và có thể kế thừa từ cấp trên.
    \item \textbf{AI không tách rời nghiệp vụ:} OCR, RAG, tóm tắt, phân loại và cảnh báo thiếu minh chứng nằm trong workflow xét duyệt.
    \item \textbf{Chủ động giúp sinh viên hoàn thiện tiêu chí:} hệ thống không chỉ báo ``thiếu'' mà còn gợi ý hoạt động đang mở từ nguồn Hội Sinh viên/trường/khoa/CLB để sinh viên kịp bổ sung trước deadline.
    \item \textbf{Tích lũy thành tích từ sớm:} sinh viên không cần chờ tới mùa xét mới gom hồ sơ; kho thành tích giúp lưu minh chứng liên tục, còn hệ thống tự lọc minh chứng hợp lệ theo chu kỳ xét.
    \item \textbf{Tạo giá trị sau khi xét duyệt:} CV/portfolio được tạo từ dữ liệu đã xác nhận giúp danh hiệu và minh chứng không chỉ nằm trong hồ sơ xét, mà trở thành tài sản năng lực của sinh viên.
    \item \textbf{Dễ mở rộng:} MVP demo 2 cấp khoa -{}> trường trong một cây tổ chức nhỏ, sau đó mở rộng dần lên cấp thành phố/tỉnh và cấp Trung ương.
    \item \textbf{Minh bạch và truy vết:} có audit log, lịch sử duyệt, danh hiệu đã đạt và gói hồ sơ gửi cấp trên.
\end{itemize}

\subsection{Mô hình doanh thu hoặc giá trị mang lại}

\begingroup
\small
\setlength{\tabcolsep}{4pt}
\renewcommand{\arraystretch}{1.25}

\begin{longtable}{|>{\raggedright\arraybackslash}p{0.25\textwidth}|
                  >{\raggedright\arraybackslash}p{0.43\textwidth}|
                  >{\raggedright\arraybackslash}p{0.24\textwidth}|}
\hline
\multicolumn{1}{|>{\centering\arraybackslash}p{0.25\textwidth}|}{\textbf{Mô hình}} &
\multicolumn{1}{>{\centering\arraybackslash}p{0.43\textwidth}|}{\textbf{Mô tả}} &
\multicolumn{1}{>{\centering\arraybackslash}p{0.24\textwidth}|}{\textbf{Phù hợp giai đoạn nào}} \\ \hline
\endfirsthead

\hline
\multicolumn{1}{|>{\centering\arraybackslash}p{0.25\textwidth}|}{\textbf{Mô hình}} &
\multicolumn{1}{>{\centering\arraybackslash}p{0.43\textwidth}|}{\textbf{Mô tả}} &
\multicolumn{1}{>{\centering\arraybackslash}p{0.24\textwidth}|}{\textbf{Phù hợp giai đoạn nào}} \\ \hline
\endhead

Thí điểm miễn phí/đồng hành với Hội Sinh viên &
Cung cấp bản thử nghiệm cho một khoa/trường hoặc một mùa xét, có đầu mối Hội Sinh viên/Đoàn - Hội cùng cấu hình tiêu chí, thử quy trình và phản hồi nghiệp vụ &
0-3 tháng sau cuộc thi \\ \hline

Thu phí theo đơn vị triển khai &
Trường/khoa/tỉnh thành trả phí theo gói cấu hình, tài khoản cán bộ, dashboard, lưu trữ, báo cáo &
Sau khi MVP ổn định \\ \hline

Thu phí theo hồ sơ/API usage &
Tính theo số hồ sơ xử lý, lượt OCR, lượt hỏi RAG, dung lượng lưu trữ &
Khi mở rộng nhiều đơn vị \\ \hline

Gói triển khai toàn hệ thống &
Cung cấp hạ tầng, tài liệu, đào tạo, hỗ trợ vận hành và quy trình onboarding cho nhiều cấp Hội Sinh viên &
Giai đoạn triển khai rộng \\ \hline

Giá trị xã hội &
Tăng tỷ lệ sinh viên theo đuổi danh hiệu, giảm tải cán bộ, minh bạch hóa phong trào &
Xuyên suốt \\ \hline

\end{longtable}
\endgroup